\chapter*{Appendix: Methodology and the Companion Repository}
\addcontentsline{toc}{chapter}{Appendix: Methodology and the Companion Repository}
\markboth{Appendix: Methodology}{}

% Long monospace file names cannot hyphenate on an A5 measure. Give TeX room
% to stretch rather than overrun, as the glossary does for its compounds.
\begingroup
\emergencystretch=3em
\hbadness=6000
\sloppy

The preface promised that where a number in these pages is mine, the
script that produced it ships with the book, and that where it belongs to
somebody else, the citation says whose. This appendix is what makes that
promise checkable. It describes what the companion repository holds and which
part of it produces which chapter's numbers. The repository is at \repourl.

\section*{How the Numbers Were Made}

The companion repository is a working file cabinet, not a software product. It
holds the data the book runs on, the scripts that turn that data into the
numbers printed in the chapters, the scripts that draw the figures, and three
checkers that police the manuscript itself. Those live under \texttt{analysis/data},
\texttt{analysis/checks} and \texttt{analysis/figures}, with two larger pieces
beside them: the Chapter~5 backtest kit, and the historical calculator that
Chapter~9 asks you to run on your own mix of its two machines. The repository outlives this
printing, so every edition of the book is published alongside a matching tag.
The edition's tag is listed on the copyright page. Check out that tag to reproduce the version used for this edition. % src: the printed edition and its tag are cut together; the tag name is set at press and printed in the front matter

The division of labor is deliberate. A fetcher goes out to a public source and
writes its result into \texttt{analysis/data}; there are ten. An analysis
script reads those files and prints to the console every number that a source comment in its chapter attributes to it, and draws nothing. These are the twenty-seven R
scripts and seven Python scripts in \texttt{analysis/checks}, which is a directory named for its purpose rather than for most of its contents. A figure script reads the same
files, or a small summary table an analysis script wrote, and draws exactly one
figure. Several figure scripts also print the values their caption quotes,
which is why none of Chapters~1 to 4 needs more than three analysis scripts of its own.
Chapter~2 never touches \texttt{analysis/checks}, Chapter~3 has one of its
own, and Chapter~1 has three: for the network's power, for the listed fleets
and the other checks on its efficiency range, and for when its hashrate peaked. Chapter~4
borrows two of them. Chapter~1 also borrows Chapter~6's ledger for the
delivered price, and Chapter~7 reaches into the directory only for the one
number it borrows from Chapter~6. % src: analysis/data/fetch_*.py and fetch_*.R (9) and the readability poller analysis/data/ercot_api_readability/poll.py (10); python3 analysis/checks/check_src_comments.py --count, analysis/checks inventory: 27 .R analysis scripts (ch9_tuning.R and ch9_hashprice.R among them, libraries the ledger scripts load, ch8_demand.R, which every chapter 8 demand script sources, official_4cp.R, the library every reader of the official 4CP intervals loads, and oncor_distribution.R, the distribution tariff that ch6_btm_ledger.R, ch9_optimal_stack.R, ch10_hashprice_sensitivity.R and ch9_calculator_export.R bill from; 26 until 2026-10-03), 7 .py analysis scripts (ch5_price_autocorrelation.py, ch5_adder_example.py, ch5_stakes.py, ch5_posting_latency.py, ch5_api_readability.py, ch5_vfloor.py, ch8_postage_stamp.py) and 3 .py checkers (check_citations.py, check_glossary_trace.py, check_src_comments.py); no script in analysis/checks calls ggsave, png() or pdf(); every analysis/figures/*.R writes one figure as .png and .pdf; same run, analysis/checks scripts the source notes name: chapter 1 ch1_network_energy.R, ch1_fleet_efficiency.R and ch1_hashrate_peak.R, and chapter 6's ch6_btm_ledger.R (chapter-1.tex:55, the delivered price), chapter 3 ch3_typical_hour.R, chapter 4 two of chapter 1's, ch1_network_energy.R and ch1_hashrate_peak.R (chapter-4.tex:24 and :42-43), chapter 7 ch6_contract_pnl.R (chapter-7.tex:91), chapter 2 none

Every quantitative
claim in the chapters carries a comment in the \LaTeX{} source, written
\texttt{\% src:}, sitting immediately after the sentence it belongs to. There
are 592 of them across the chapters and the preface, and 17 more in this
appendix. % src: python3 analysis/checks/check_src_comments.py --count: 609 src comments across book/chapters/*.tex: 591 in chapters 1-10 + one in the preface = 592, the figure printed here, and 17 in this appendix; glossary.tex and the roadmap carry none; counted as comments, not string matches (a comment whose text, after the first unescaped percent sign on its line, starts with src:): a string match also finds the marker this appendix prints and a second marker inside one chapter 5 comment (chapter-5.tex:127), 611 in all, and the run lists both by line; fig-src and tab-src notes (26 and 7) are not in these totals
Most name a bibliography key, usually with a section or page
number, or the path of the script that computed the number; a few name the
committed data file directly, and a few show the arithmetic inline where the
step is one line long and a script would be theater. They live in the source
and never appear on the printed page, so tracing a number means opening the
chapter's \texttt{.tex} file rather than hunting through these pages. The
glossary is the exception: its entries restate numbers the chapters already
source, and carry no comments of their own. One of the three checkers,
\texttt{check\_glossary\_trace.py}, finds the numbers the glossary prints and
the chapter line each restates, and fails on any it cannot place or whose
chapter line has changed since it was last checked by hand. % src: python3 analysis/checks/check_glossary_trace.py: 635 quantitative tokens in 184 glossary entries, 608 traced in context to a chapter line, 27 pinned by the 27 rows of analysis/checks/glossary_trace_exceptions.csv (roundings, counts and positions the chapter gives by naming things in order, complements, and real homes the matcher passes over for a coincidence, each row checked against verbatim text of its chapter line), 0 value-only, 0 untraced; every row is held against the reviewed trace map analysis/data/glossary_trace.csv, whose best hits were read by hand

The mix of source types differs from chapter to chapter. Chapters~1 through
3 carry 107 source comments between them, and only 18 name a script or
a committed data file, because most of their numbers belong to other people.
Chapter~9 names a script in 57 of its 69, and a script or a committed data file
in 60. % src: python3 analysis/checks/check_src_comments.py --count: chapter-1/2/3 27 + 40 + 40 = 107 src comments, eighteen naming a script or an analysis/data file (chapter 1 seventeen, chapter 3 one, chapter 2 none); chapter-9.tex 69 src comments, 57 naming a .R or .py script, 60 naming a script or an analysis/data file

% The appendix is unnumbered, so its table would continue chapter 10's count.
\renewcommand{\thetable}{A.\arabic{table}}\setcounter{table}{0}
\begin{table}[htbp]
  \centering
  \scriptsize
  \setlength{\tabcolsep}{3.5pt}
  \begin{tabular}{@{}p{11mm}p{88mm}@{}}
    \hline
    \textbf{Ch.} & \textbf{What produces its numbers and figures} \\
    \hline
    1 & \texttt{ch1\_network\_energy.R} (the network's power and energy across
        the range of its fleet efficiency), \texttt{ch1\_fleet\_efficiency.R}
        (the checks on that range: the efficiency the listed fleets
        publish, the share of revenue each end of it would spend on power, and
        Cambridge's figures, quoted) and
        \texttt{ch1\_hashrate\_peak.R} (when its hashrate peaked); figures from
        \texttt{fig\_ch1\_network\_hr\_diff.R},
        \texttt{fig\_ch1\_asic\_efficiency.R} (on a hand transcription of
        Bitmain's own documents), \texttt{fig\_ch1\_global\_generation.R}. The
        listed fleets are hand transcribed from the companies' filings and
        releases, their power bills from a fund manager's compilation of
        those filings. The other numbers are citations, and
        the machine specifications are the manufacturers' own. \\
    2 & \texttt{fig\_ch2\_uri\_frequency.R}, a schematic reconstruction from the
        FERC/NERC staff report's anchor values. The numbers are citations. \\
    3 & \texttt{fig\_ch3\_lz\_map.R}, and \texttt{ch3\_typical\_hour.R} for
        what a typical hour clears at. The other numbers are protocol, tariff
        and regulatory citations. \\
    4 & \texttt{fig\_ch4\_hashprice\_cycle.R} (bitcoin price, network hashrate,
        hashprice) and \texttt{fig\_ch4\_sensitivity.R} (margin against
        hashprice for the two standing machines). \\
    5 & \texttt{ch5\_oracle\_check.R}, a simulator-free check on how often
        the run-or-pause call is in doubt; \texttt{ch5\_price\_autocorrelation.py},
        \texttt{ch5\_adder\_example.py} and \texttt{ch5\_stakes.py} for the
        price statistics, and \texttt{ch5\_posting\_latency.py} for when ERCOT
        posts each print and \texttt{ch5\_api\_readability.py} for when it can be read; \texttt{ch5\_vfloor.py} for the fleet data below each
        machine's best-efficiency point; \texttt{ch5-backtest/} for the control-loop study itself
        (\texttt{papo\_sim.py}, eight study drivers, an optimizer check and a
        database builder); five
        \texttt{fig\_ch5\_*.R} scripts for the efficiency curves, the optimal
        power target, the hold and price deadband, the backstop, and the oracle
        benchmark. \\
    6 & \texttt{ch6\_btm\_ledger.R} (what a 100 MW mine pays on top of
        energy), \texttt{ch6\_contract\_pnl.R} (cost by contract shape),
        \texttt{ch6\_basis.R} (basis), \texttt{ch6\_strike\_events.R} (the
        strike's pauses), \texttt{sell\_vs\_mine.R} (the owner-generator's
        choice), \texttt{oncor\_distribution.R} (the distribution bill, shared with Chapters~9 and~10); four \texttt{fig\_ch6\_*.R} scripts. \\
    7 & \texttt{fig\_ch7\_as\_history.R} and
        \texttt{fig\_ch7\_standing\_revenue.R}. Both print the clearing prices
        and annual standing revenues the text quotes. \\
    8 & \texttt{ch8\_postage\_stamp.py} sets the postage stamp rate against
        the cost of service it recovers; \texttt{ch8\_demand.R} assembles the
        complete demand proxy from ERCOT fuel generation, without gap-fill; \texttt{official\_4cp.R} reads ERCOT's official
        intervals and derives each one's start from its label, for this
        chapter, Chapters~9 and 10; \texttt{ch8\_4cp\_peaks.R} reconstructs the monthly peak intervals;
        \texttt{ch8\_4cp\_prices.R}, \texttt{ch8\_4cp\_rules.R} and
        \texttt{ch8\_12cp\_haystack.R} price them, test forecast-free dodging
        rules, and count how thick the near-peak haystack gets under twelve
        peaks; \texttt{ch8\_ers\_blocks.R} prices ERS. Figures from \texttt{fig\_ch8\_4cp\_calendar.R} and
        \texttt{fig\_ch8\_4cp\_window.R}. \\
    9 & \texttt{ch9\_flexibility\_stack.R} and \texttt{ch9\_optimal\_stack.R}
        carry the ledger; \texttt{ch9\_tuning.R} is the power-target layer that they,
        the cross-checks, the calculator export and Chapter~10 all load, and
        \texttt{ch9\_hashprice.R} is the daily hashprice every ledger
        script reads, with the calendar checked so that no day is dropped
        unseen (\texttt{ch9\_vs\_papo.R} excepted: it reproduces an
        external run on that run's own 362 days);
        \texttt{ch9\_control\_friction.R} and \texttt{ch9\_vs\_papo.R} are the
        cross-checks; \texttt{ch9\_calculator\_export.R} builds the
        calculator's payload. Four \texttt{fig\_ch9\_*.R} scripts. \\
    10 & \texttt{ch10\_subsidy\_arithmetic.R} (the halving schedule and the fee
        share of the security budget) and \texttt{ch10\_hashprice\_sensitivity.R}
        (Chapter~9's reference site rerun at lower hashprice), with
        \texttt{fig\_ch10\_security\_budget.R} and
        \texttt{fig\_ch10\_hashprice\_sensitivity.R}. \\
    \hline
  \end{tabular}
  \caption{What produces each chapter's numbers and figures. Scripts live under
  \texttt{analysis/checks} and \texttt{analysis/figures} in the companion
  repository, and Chapter~5's study under \texttt{analysis/ch5-backtest}.} % src: file listing of analysis/checks, analysis/figures and analysis/ch5-backtest (papo_sim.py; drivers mixed_basis_study.py, naive_gap.py, baselines_study.py, capture_anatomy.py, trend_followup.py, fade_study.py, multizone_study.py, reactive_papo_trend.py; optimizer_bruteforce_check.py; build_prices_db.py); header comment of each script names its chapter and its output
  \label{tab:appendix-scripts}
\end{table}

ERCOT settlement point prices for all
eight load zones, day-ahead and real-time back to 2021, and for two hubs in real time, come from ERCOT's own
public archive through \texttt{fetch\_ercot\_rtm15\_spp\_history.py}, which
requires no login credentials or API key: these are files ERCOT posts on the
open web, and that one script writes all of these price files. Ancillary-service clearing
prices come from the same archive through
\texttt{fetch\_ercot\_as\_mcpc\_history.py}. Fifteen-minute generation by fuel comes from
ERCOT's Fuel Mix Report through \texttt{fetch\_ercot\_fuel\_mix.py}.
Chapter~8 uses the sum, including negative battery charging, as a demand proxy:
net DC-tie imports and explicit Block Load Transfer adjustments are omitted.
The early workbooks include storage-load adjustments inside ``Other''; later
workbooks report them separately. The sum includes them in either case.
Every operating day is complete, with no gap-fill.
The latest months are still ERCOT's initial settlement and are not presented as final. The repository also keeps an hourly Native Load snapshot, which no printed number uses.
The moment ERCOT posts each five-minute price comes from its public report
listing through \texttt{fetch\_ercot\_np6788\_posting.py}. None of these
needs an account or key. Three scripts do need a free ERCOT API
Explorer account and a connection from inside the United States. The poller in
\texttt{analysis/data/ercot\_api\_readability} timed when each price became
readable; its ten days of polls ship, so checking those numbers needs no
account. \texttt{fetch\_lz\_west\_ercot.py} and the backtest kit's own ingester
are on the path to no printed number; they are there to extend the series past
what ships. % src: analysis/data/fetch_ercot_rtm15_spp_history.py docstring ("No account, no API key, no credentials: these are public files"; MIS reportTypeId 13061 and 13060, two outputs; with --hubs, only the third, ercot_rtm15_spp_hubs_2021on.csv.gz, HB_NORTH and HB_WEST from 13061, read by analysis/checks/ch6_basis.R); fetch_ercot_as_mcpc_history.py docstring (MIS reportTypeId 13091, NP4-181-ER Historical DAM Clearing Prices for Capacity, one DAMASMCPC_YYYY.zip a year; no account, no key); fetch_ercot_native_load.R header ("Credential-free"), fetch_ercot_fuel_mix.py docstring ("No account, no key, no credentials") and fetch_ercot_np6788_posting.py docstring ("Credential-free", MIS reportTypeId 12300 listing); fetch_lz_west_ercot.py and analysis/ch5-backtest/ingest/ercot_ingest.py read ERCOT_USERNAME / ERCOT_PASSWORD / ERCOT_SUBSCRIPTION_KEY

Every ERCOT series in the repository comes from ERCOT itself. The
five-minute prints, the 15-minute settlement point prices and the
ancillary-service and ERS clearing prices are pulled from its Market Information
System, and the annual aggregates that Chapters~4 and 6 use as cross-checks are
rebuilt from those files by a script rather than carried as
opaque tables. Hashprice,
network hashrate and difficulty come from the Braiins Insights public API and
from blockchain.info's public charts. Plant locations are EIA-860M and the Texas
outline is a Census boundary file, both public domain. The efficiency curves come
from the fleet dataset, aggregated operating data from real machines running
Braiins~OS rather than measurements in a laboratory, and the chapters say so every time they use them. % src: analysis/data/README.md provenance table; analysis/data/build_lz_west_annual_means.R rebuilds both annual-means tables from ercot_spp_rtm5_lzwest.csv.gz and ercot_dam_spp_loadzones_2021on.csv.gz, reproducing the 2025 real-time figures exactly; Braiins Insights; EIA-860M June 2026; Census cb\_2023\_us\_state\_20m; ercot_dam_as_mcpc.csv.gz from MIS 13091 (fetch_ercot_as_mcpc_history.py), ercot_ers_procurement.csv from MIS 11465, NP3-144-M (fetch_ercot_ers_procurement.py, credential-free), analysis/data/README.md fetchers table; ercot_fuel_mix_15min_2019on.csv.gz from the Fuel Mix Report, 2019-2021 WSL adjustments included in Other according to each workbook Summary note, FINAL through July 2026 and INITIAL August-September 2026

Braiins~OS sends operating data about the machines it runs to Braiins, which uses
it to develop the firmware, as its privacy policy discloses \autocite{braiins2020osprivacy}.
This book uses only the fleet dataset: one average per machine model and power
target for the second quarter of 2026, committed to the companion repository,
with no single customer's or machine's figures in it. % src: braiins2020osprivacy ("When you use your virtual currency mining device with Braiins OS operating system to mine virtual currencies, we may collect and process telemetric data related to your mining activities"; "We collect, store and process the above-mentioned analytical data to develop and test our Braiins OS products"); analysis/data/miner_efficiency_telemetry_2026q2.csv, columns model, power_w, avg_jth and hashrate_ths, 35 rows over two models (S19j Pro, S21 XP), no customer or machine identifier; wording proposed by Kristian Csepcsar (Zulip, Vit's Publications > Book, 2026-10-13) after a review comment that "telemetry" had drawn community criticism

Some inputs are committed snapshots rather than live calls so the results can be reproduced. \texttt{fig\_ch1\_network\_hr\_diff.R} is the
worked example. Its endpoint returns a rolling window that always ends at the
moment of the request, so a script that fetched at build time would draw a
different chart every run and the printed figure could never be checked against
anything. It therefore reads a committed file, retrieved on September 14, 2026,
and refreshes it only when you ask for that explicitly, which leaves the change
visible in the repository's history. % src: analysis/figures/fig_ch1_network_hr_diff.R header and the REFRESH=1 branch; snapshot analysis/data/braiins_hashrate_difficulty_all.json, 104 points to 2026-09-14

Before the Chapter~5 studies will run, one file has to be built. The price
database they read is about 230 megabytes, too large to ship, so
\texttt{build\_prices\_db.py} reconstitutes it in well under a minute from
files that are committed: ERCOT's own settlement archive for the 15-minute
and day-ahead prices, and the five-minute prints. It is the first command in
the kit's instructions and the only setup step in the repository. % src: analysis/ch5-backtest/build_prices_db.py console (2,250,452 rows from ercot_rtm15_spp_loadzones_2021on.csv.gz, ercot_dam_spp_loadzones_2021on.csv.gz and ercot_spp_rtm5_lzwest.csv.gz; 227 MB; about 20 s, 20.3 s in a rebuild on 2026-09-25); analysis/ch5-backtest/README.md "Run it" section; analysis/data/prices.db is listed in .gitignore

The analysis has several limitations.

The two backtests are simulations, not accounts: Chapter~5's study and
Chapter~9's ledger run today's two machines through past ERCOT prices, and no
real site reported any of their figures. The oracle, the perfect-foresight twin
and the hindsight ceiling see the outcome first, so each is
printed as a ceiling, never as a strategy. Chapter~5's rules would have cost
Chapter~9's two fleets 0.15~percent of their frictionless mining margin before grid charges in
2025, with no charge for moving a target. Chapter~9's ledger covers 2021 to
2025, Chapter~5's study runs to June 2026, and 2026 elsewhere is its first
eight months, printed as such. % src: chapter-9.tex:11 (the fleet a control held fixed across the five years; what today's fleet would have earned, not what a site of the time did); chapter-5.tex:133-134 (the oracle: perfect foresight, a ceiling, not a strategy); chapter-9.tex:40 and :140 (the perfect-foresight twin; the best annual choice in hindsight printed as a ceiling, never as a plan, and never called a strategy); analysis/checks/ch9_control_friction.R console, rerun 2026-09-27 (2025 only, each fleet run alone at 74.10 MW: frictionless $53,741,258, chapter 5 rules $53,660,694, -$80,564 = 0.15%; moves counted, never charged), as chapter-9.tex:79-80; analysis/checks/ch9_optimal_stack.R console header (ledger years 2021-2025, 2021 from 2021-01-02), chapter-9.tex:41; analysis/ch5-backtest/README.md (TEST-2025-26 ends 2026-06-22 Central), chapter-5.tex:167; 2026 elsewhere: chapter-5.tex:61 src (2026-01-01 to 08-29), chapter-9.tex:70 and :75 (the eight months of 2026), chapter-10.tex:43, :49, :74 and :96 (the first eight months, 2026 to date)

The Chapter~3 map is the weakest artifact in the book and should be read as an
orientation sketch. Load zones are electrical constructs, sets of buses, and
ERCOT publishes no official geographic data for them; the polygons drawn here
come from the only public vector source, a third-party re-host of a tracing of
ERCOT's published map, with no license attached to it. The four
non-opt-in zones have no public polygons at all and appear only in the caption.
The dots are not nodes either: ERCOT withholds true bus coordinates as secured
infrastructure information, so the map shows power plants from EIA-860M
standing in for the generation side. % src: analysis/figures/fig_ch3_lz_map.R provenance header, layers 2 and 3 (ICF ArcGIS layer, unlicensed re-host, geometry traced from ERCOT's published map; NP4-160-SG; LZ_AEN, LZ_CPS, LZ_LCRA and LZ_RAYBN have no public polygons; node dots approximated by 988 EIA-860M plants)

ERCOT's 4CP workbooks since 2008, the NP9-83-M files among them, leave one
convention unstated, and every price and rule window that Chapters~8 and 9
pair with a peak depends on it. A time such as 17:00 names the end of the
15-minute settlement interval, so the official interval labeled 17:00 ran
from 4:45 to 5~p.m. ERCOT's archived workbooks for 2003 to 2006, and its
preliminary ones for 2007, say so in words, ``at interval ending'' beside the
same bare time; its 2016 report to the commission does too; the protocols define
every settlement interval by the time it ends; and a wires utility read
ERCOT's 2024 times the same way in a 2025 rate filing
\autocite{ercot4cphistzip,ercot4cp2016puct,ercot2026protocols,aeptexas58315}.
The committed demand series agrees. On 23 of the 28 official days of 2019 to
2025, the quarter-hour ending at the label carried more demand than the one
starting at it. The repository keeps ERCOT's labels verbatim and derives each
interval's start in one place, \texttt{official\_4cp.R}, which every script
that pairs a peak with a price or a rule window reads. % src: ercot4cphistzip (2003-2006 workbooks and the June-August 2007 preliminary workbooks: at interval ending; 2006/FILING_20064CP.xls JUNE sheet, Peak Demand on 06/13/2006 at interval ending 16:15, beside the summary header 6/13/2006 16:15; the 2008-2020 workbooks and ercotnp983m 2022-2025 print bare labels only); ercot4cp2016puct p. 2 (Interval Ending (IE) 16:00; its workbook labels it 7/14/2016 16:00); ercot2026protocols sec. 10.9.1(2)(h) and 10.9.2(1)(h); aeptexas58315 p. 2 (2024 matrix peaks at the 15-minute interval ending times); analysis/checks/ch8_4cp_peaks.R console, label convention check (ending beats starting on 27 of 28; the day's peak ends at the label on 24 and starts at it on 1); analysis/checks/official_4cp.R

Neither backtest engine charges anything for moving a power target. Chapter~5's
and Chapter~9's both apply a new target within the interval, at the full
efficiency of the curve. A machine that has no saved profile for a target
spends a settling period tuning to it, so a real fleet can earn less than the
engines show, by an amount that depends on the firmware and on how often the
loop lands on a target it has not held before. In Chapter~5's kit that is once
per corner point of the curve, four targets on the S21~XP and eight on the
S19j~Pro; Chapter~9's engine chooses from a 100~W grid and is not bound to
those points. Nothing in the repository measures the settling cost
itself. % src: analysis/ch5-backtest/papo_sim.py simulate() (applied = ideal, no per-move charge); analysis/ch5-backtest/optimizer_bruteforce_check.py console, "targets the loop applies" block (4 of 4 and 8 of 8 hull points, no other target after each run's first move), as chapter-5.tex, Three Facts, prints them; analysis/checks/ch9_control_friction.R (moves counted, never charged); analysis/checks/ch9_tuning.R ch9_curve (J/TH interpolated on a 100 W grid, any grid point selectable) takes the best operating point in every interval, with no charge for changing it; the settling period as chapter-5.tex, Three Facts (braiinsacademyautotuning)
Nor does either engine charge a paused machine for what it still draws. In
Chapter~5's kit, each percent of stock power drawn through every pause, billed
at the settlement point price, would take about 1.1 percent off the full loop's gain
over stock power on the S19j~Pro and 0.4 percent on the S21~XP, out of sample;
the repository holds no measurement of that draw. % src: analysis/ch5-backtest/baselines_study.py residual block (papo_sim.py simulate(paused_draw=0.01), TEST-2025-26 LZ_WEST: -1.10% (S19j Pro) and -0.43% (S21 XP) of the full loop's gain over S0 per 1% of stock draw, billed at the settlement price; 2%: -2.21% and -0.87%; TUNE-2024: -0.81% and -0.29% per 1%); papo_sim.py simulate() books a paused interval at 0 by default; analysis/checks/ch9_tuning.R prof_tuned and prof_stock are 0 when a fleet is off

Two bodies of material in the repository are Braiins' own: the fleet
dataset behind the efficiency curves, and the Chapter~5 backtest kit. Both are
published here for the first time. Where a chapter rests on either, it says so
in the text and not only in a source comment. % src: analysis/data/README.md provenance table (fleet telemetry row) and analysis/ch5-backtest/README.md; chapter-5.tex provenance paragraph

\section*{On the Use of AI}

I used Claude Code, an AI coding agent, throughout this project: to find
candidate sources, to assemble the bibliography, to write the analysis and
figure scripts, to audit that bibliography afterwards, and as a drafting
assistant on the prose.

\endgroup
